A player keeps a Perch account in \fact{wallet-name} and wants Guess the Number,
a Soroban game that takes \fact{game-wrong-guess-xlm}~XLM per wrong guess and
pays out its balance on a right one, to play without a passkey prompt per turn.
The owner can install a \emph{grant} naming the contract and functions an app may
call, the signer that may act, how much it may spend over time and when it
ends. It should not let the game call the account's own
entry points or its recovery controller. We follow three grants that the
reference run of \S\ref{sec:evaluation} installs. Grant~1 lets an app transfer
at most \fact{grant1-cap-xlm}~XLM in any \fact{grant1-period-ledgers}-ledger
window (\fact{grant1-period-minutes} at the nominal
\fact{nominal-ledger-seconds}-second ledger): a
\fact{grant1-spend-applied-xlm}-XLM transfer is applied, and a further
\fact{grant1-spend-refused-xlm}~XLM in that window is refused. Grant~2 gives that
app a fresh key after recovery. Grant~3 lets the game call \code{guess} and
transfer at most \fact{grant3-cap-xlm}~XLM per \fact{grant3-period-ledgers}
ledgers, so two lost guesses fit and a third is refused. Two further
requirements follow: a grant derived from another must allow nothing its parent
does not, and a completed recovery must end earlier grants before their storage
is collected.

\paragraph{Problem and approach.} A grant needs one meaning in three places:
across the nested calls of one authorization, which Soroban asks the account to
decide one at a time; under rule selection, which whoever builds the transaction
performs; and across recovery, which replaces the credentials a grant was issued
under. A spending cap also needs state. Perch makes the authorization decision a
pure state transition and separates what is proved about it from the
obligations of the account that invokes it. An evaluator linked into the account
maps the selected rules, their spending records, the calls, the claimed signers,
the ledger and a grant generation to a refusal or the next spending records,
before any signer is authenticated (Fig.~\ref{fig:fig-parts}).

\perchfigure{fig-parts}

\paragraph{Contributions.}
\begin{enumerate}[nosep,leftmargin=*]
\item A Soroban smart-account design in which authorization is a pure state
  transition linked into the account. It evaluates the selected rules against
  account-held spending records and a grant generation before authenticating
  any signer. Grant lowering excludes the account, its adopted recovery
  controller, and the zero-knowledge adapter and pool of its recovery
  configuration (\S\S\ref{sec:grants}--\ref{sec:authorization}).
\item A grant generation that a completed recovery advances. Every earlier
  grant, and every grant derived from one, stops authorizing in the completion
  transaction, and a credential index bars their keys from later documents
  (\S\ref{sec:recovery}).
\item Lean theorems over the Aeneas translation of the Rust evaluator:
  equivalence with a model, acceptance implies permission, attenuation and trace
  inclusion, window-cap safety, and containment. They come with an explicit
  trust base, and \fact{lean-doc-traces} Lean-generated traces connect them to
  the account (\S\ref{sec:proofs}).
\item A source-cited comparison with oz-policy-builder in architecture,
  expressiveness, binding and assurance, with one task run in both designs in
  one harness (\S\ref{sec:compare}), and an evaluation on a local Soroban host
  and on testnet through a wallet preview (\S\ref{sec:evaluation}). Comparative
  cost measurement beyond that task is proposed, not reported.
\end{enumerate}

\paragraph{Research questions.}
\textbf{RQ1} (enforcement): does the built account enforce grants, caps and
recovery invalidation as the design states? \textbf{RQ2} (assurance): which
parts of a decision are proved, which tested, and which trusted? \textbf{RQ3}
(cost): what does a decision cost on chain, for Perch alone? \textbf{RQ4}
(comparison): how does a linked evaluator differ from an attached shared
interpreter in expressiveness, binding, state and assurance? \S\ref{sec:evaluation}
answers RQ1 and RQ3, \S\ref{sec:proofs} RQ2 and \S\ref{sec:compare} RQ4.

The design is an open pull-request stack, previewed on Stellar testnet
(\S\ref{sec:status}). After background, threat model and status
(\S\S\ref{sec:background}--\ref{sec:status}), we present the design
(\S\S\ref{sec:grants}--\ref{sec:recovery}), the comparison, formal results and
evaluation (\S\S\ref{sec:compare}--\ref{sec:evaluation}), then related work,
limitations and future work.
